\documentclass[../../../include/open-logic-section]{subfiles}

\begin{document}

\olfileid{sth}{choice}{vitali}
\olsection{পরিশিষ্ট: ভিতালির কূটাভাস}

বানাখ--তার্স্কি নির্মাণ গ্রহণযোগ্য কি না
সত্যিই বুঝতে হলে তার \emph{প্রমাণ} দেখা উচিত।
দুর্ভাগ্য, তাতে এখানে যতটা দেওয়া সম্ভব,
তার চেয়ে অনেক বেশি বীজগণিত লাগবে। তবে
নিচের ফলটিতে মন দিয়ে বানাখ--তার্স্কির
প্রমাণের ধারণা বোঝাতে কয়েকটি কথা বলা যায়।\footnote{
আরও পূর্ণ আলোচনার জন্য দেখো \cite{Weston2003}
বা \cite{Wagon2016}।}

\begin{thm}[ভিতালির কূটাভাস ($\ZFC$-এ)]\ollabel{vitaliparadox}
ধনাত্মক ব্যাসার্ধের যেকোনো বৃত্তকে তালিকায়নযোগ্য সংখ্যক খণ্ডে
ভাগ করে, সেগুলি ঘুরিয়ে ও স্থানান্তর করে
আবার সাজালে সেই বৃত্তের দুটি প্রতিলিপি
গঠন করা যায়।
\end{thm}

ভিতালির কূটাভাসের প্রমাণ বানাখ--তার্স্কির
প্রমাণের চেয়ে অনেক সহজ। একে ``ভিতালির
কূটাভাস'' বলেছি, কারণ ভিতালির
\citeyear{Vitali1905}-এর অপরিমাপযোগ্য
সেট-নির্মাণ থেকে এটি মেলে। তবে দুটি
প্রমাণের সেটতাত্ত্বিক অংশ খুব কাছাকাছি।
মূল পার্থক্য, বানাখ--তার্স্কিতে \emph{সসীম}
বিভাজন, ভিতালির কূটাভাসে \emph{তালিকায়নযোগ্য
অসীম} বিভাজন। \citet{Weston2003}-এর
কথায়, ভিতালির কূটাভাস ``বানাখ--তার্স্কির
মতো অতটা চমকপ্রদ নয়, কিন্তু এটি দেখায় যে
`সরল' অবস্থাতেও জ্যামিতিক কূটাভাস ঘটতে পারে।''

ভিতালির কূটাভাস দ্বিমাত্রিক একটি চিত্র,
বৃত্ত নিয়ে। তাই সমতল $\Real^2$-এ কাজ করব।
% BN-SRC-472: rational multiples of a full turn, not rational radian
% values, are closed under composition and inversion modulo a full turn.
$\rotationsgroup$ হলো মূলবিন্দুকে কেন্দ্র করে
বিন্দুগুলিকে $[0,2\pi)$-এর মধ্যে পূর্ণ আবর্তনের
\emph{মূলদ ভগ্নাংশের} সমান কোণে ঘড়ির কাঁটার
দিকে ঘোরানোর সমষ্টি। $\rotationsgroup$ সম্বন্ধে
কিছু বীজগাণিতিক তথ্য দেখি (বক্তব্যটি না বুঝলে
প্রমাণে তার অর্থ স্পষ্ট হবে):

\begin{lem}\ollabel{rotationsgroupabelian}
অপেক্ষকের সংযোজন ক্রিয়ায়
$\rotationsgroup$ একটি আবেলীয় {দল}।
\end{lem}

\begin{proof}
$0_{\rotationsgroup}$ দিয়ে $0$ রেডিয়ানের
ঘূর্ণন বোঝাই। এটি $\rotationsgroup$-এর
অভেদ সদস্য, কারণ যেকোনো
$\rho\in\rotationsgroup$-র জন্য
$\comp{0_{\rotationsgroup}}{\rho}=
\comp{\rho}{0_{\rotationsgroup}}=\rho$।

প্রত্যেক সদস্যের বিপরীত আছে।
$\rho\in\rotationsgroup$ যদি $r$ রেডিয়ান
ঘোরায়, তবে অশূন্য ঘূর্ণনের বিপরীত
$\rho^{-1}\in\rotationsgroup$
$2\pi-r$ রেডিয়ান ঘোরায়; শূন্য ঘূর্ণনের
বিপরীত সেটি নিজেই। ফলে
$\rho\circ\rho^{-1}=0_\rotationsgroup$।

সংযোজনে বন্ধনী বদলেও একই ফল: যেকোনো
$\rho,\sigma,\tau\in\rotationsgroup$-র
জন্য $\comp{\rho}{(\comp{\sigma}{\tau})}=
\comp{(\comp{\rho}{\sigma})}{\tau}$।

সংযোজন বিনিময়ী: যেকোনো
$\rho,\sigma\in\rotationsgroup$-র জন্য
$\comp{\rho}{\sigma}=\comp{\sigma}{\rho}$।
\end{proof}

বস্তুত $\rotationsgroup$-কে দুই ভাগে
ভাগ করেও প্রত্যেক ভাগ দিয়ে পুরো দলটি
উদ্ধার করা যায়:

\begin{lem}\ollabel{disjointgroup}
$\rotationsgroup$-এর দুটি বিযুক্ত ভাগ
$\rotationsgroup_{1}$ ও
$\rotationsgroup_{2}$ আছে; প্রত্যেক ভাগই
উপযুক্ত সংযোজনে পুরো
$\rotationsgroup$-কে উৎপন্ন করে।
\end{lem}

\begin{proof}
$\rotationsgroup_{1}$-এ $[0,\pi)$-এর
সেই ঘূর্ণনগুলি নিই, যাদের কোণ পূর্ণ আবর্তনের
মূলদ ভগ্নাংশ; আর
$\rotationsgroup_{2}=\rotationsgroup
\setminus\rotationsgroup_1$। প্রাথমিক
বীজগণিত থেকে
$\Setabs{\comp{\rho}{\rho}}{\rho\in
\rotationsgroup_1}=\rotationsgroup$।
$\rotationsgroup_{2}$-এর ক্ষেত্রেও একই ফল।
\end{proof}

দলের এই তথ্য দিয়ে \olref{vitaliparadox}
প্রতিষ্ঠা করব। $\onesphere$ হলো একক বৃত্ত:
সমতলের মূলবিন্দু থেকে ঠিক~$1$ একক দূরের
বিন্দুগুলির সেট,
$\Setabs{\tuple{r,s}\in\Real^2}{\sqrt{r^2+s^2}=1}$।
$\onesphere$-কে ভাগ করতে $\onesphere$-র
উপর নিচের সম্পর্ক ধরি:
\[
	r \sim s \emph{ তখনই যখন }(\exists \rho \in \rotationsgroup)\rho(r) = s.
\]
অর্থাৎ মূলবিন্দুকে কেন্দ্র করে উপরের মূলদ-ভগ্নাংশ
ঘূর্ণনে এক বিন্দু থেকে অন্যটিতে যাওয়া যায়
তখনই $\onesphere$-র বিন্দুদুটি সম্পর্কে যুক্ত।
সহজেই দেখা যায়:

\begin{lem}
$\sim$ একটি তুল্যতা-সম্পর্ক।
\end{lem}

\begin{proof}
\olref{rotationsgroupabelian} থেকে সরাসরি।
\end{proof}

এবার চয়ন ব্যবহার করে $\onesphere$-র
$\sim$-তুল্যতা শ্রেণিগুলির প্রত্যেকটি থেকে
ঠিক একজন সদস্য নিয়ে একটি সেট $C$ নিই।
অর্থাৎ তুল্যতা শ্রেণিগুলির সেটে একটি
চয়ন-অপেক্ষক $f$ নিই,\footnote{
$\rotationsgroup$ !!{enumerable} বলে
$E$-র প্রত্যেক !!{element} !!{enumerable}।
$\onesphere$ !!{nonenumerable} বলে
\olref{vitalicover} এবং
\olref[card-arithmetic][simp]{kappaunionkappasize}
থেকে $E$-ও !!{nonenumerable}। তাই এখানে
\emph{অতালিকায়নযোগ্য} চয়ন ব্যবহৃত।}
\[
	E = \Setabs{\equivrep{r}{\sim}}{r \in \onesphere},
\]
এবং $C=\ran{f}$ ধরি। প্রত্যেক
$\rho\in\rotationsgroup$-র জন্য
$\funimage{\rho}{C}$ হলো $C$-র সব
বিন্দুকে $\rho$ ঘূর্ণনে পাঠিয়ে পাওয়া সেট।
পরের দুটি ফল বলবে, এই সেটগুলি বৃত্তটিকে
সম্পূর্ণ এবং পুনরাবৃত্তি ছাড়া ঢেকে ফেলে:

\begin{lem}\ollabel{vitalicover}
$\onesphere=\bigcup_{\rho\in\rotationsgroup}
\funimage{\rho}{C}$।
\end{lem}

\begin{proof}
$s\in\onesphere$ স্থির করি। কোনো
$r\in C$ আছে যে $r\in\equivrep{s}{\sim}$;
অর্থাৎ $r\sim s$, অর্থাৎ কোনো
$\rho\in\rotationsgroup$-র জন্য
$\rho(r)=s$।
\end{proof}

\begin{lem}\ollabel{vitalinooverlap}
$\rho_1\neq\rho_2$ হলে
$\funimage{\rho_1}{C}\cap
\funimage{\rho_2}{C}=\emptyset$।
\end{lem}

\begin{proof}
$s\in\funimage{\rho_{1}}{C}\cap
\funimage{\rho_{2}}{C}$ ধরি। তাহলে কোনো
$r_{1},r_{2}\in C$-র জন্য
$s=\rho_{1}(r_{1})=\rho_{2}(r_{2})$। তাই
$\rho^{-1}_2(\rho_1(r_1))=r_2$ এবং
$\comp{\rho_1}{\rho^{-1}_2}
\in\rotationsgroup$; সুতরাং
$r_{1}\sim r_{2}$। কিন্তু $C$ প্রত্যেক
$\sim$-শ্রেণি থেকে ঠিক একজন সদস্য বেছে
নিয়েছে। তাই $r_1=r_2$। ফলে
$s=\rho_1(r_1)=\rho_2(r_1)$,
এবং তাই $\rho_1=\rho_2$।
\end{proof}

এবার বৃত্তে দলের আগের তথ্য প্রয়োগ করি:

\begin{lem}\ollabel{pseudobanachtarski}
$\onesphere$-র দুটি বিযুক্ত ভাগ
$D_{1}$ ও $D_{2}$ আছে। $D_{1}$-কে
তালিকায়নযোগ্য সংখ্যক সেটে ভাগ করে
ঘুরিয়ে $\onesphere$-র একটি প্রতিলিপি
গড়া যায়; $D_{2}$-র ক্ষেত্রেও তাই।
\end{lem}

\begin{proof}
\olref{disjointgroup}-এর
$\rotationsgroup_{1}$ ও
$\rotationsgroup_2$ ব্যবহার করে নিই:
\begin{align*}
	D_{1} &= \bigcup_{\rho \in \rotationsgroup_1} \funimage{\rho}{C} &
	D_{2} &= \bigcup_{\rho \in \rotationsgroup_2} \funimage{\rho}{C}
\end{align*}
\olref{vitalicover} অনুযায়ী এটি
$\onesphere$-র বিভাজন; আর
\olref{vitalinooverlap} অনুযায়ী
$D_1$ ও $D_2$ বিযুক্ত। নির্মাণ থেকেই
$D_1$-কে প্রত্যেক
% BN-SRC-474: the source's R_1 is the defined rotationsgroup_1.
$\rho\in\rotationsgroup_1$-এর জন্য
$\funimage{\rho}{C}$-তে ভাগ করা যায়।
এগুলি ঘুরিয়ে পুরো $\onesphere$ পাওয়া যায়,
কারণ \olref{disjointgroup} ও
\olref{vitalicover} থেকে
$\onesphere=\bigcup_{\rho\in
\rotationsgroup}\funimage{\rho}{C}=
\bigcup_{\rho\in\rotationsgroup_1}
\funimage{(\comp{\rho}{\rho})}{C}$।
$D_2$-র ক্ষেত্রেও একই যুক্তি।
\end{proof}\noindent
এ থেকেই ভিতালির কূটাভাস মেলে।
$\onesphere$-কে তালিকায়নযোগ্য সংখ্যক
খণ্ডে (বিভিন্ন $\funimage{\rho}{C}$-তে)
ভাগ করে অনমনীয় গতিতে সাজালেই
$\onesphere$-র \emph{দুটি} প্রতিলিপি
পাওয়া যায়: $D_1$-এর খণ্ডগুলি শুধু
ঘোরাই; $D_2$-এর খণ্ডগুলি প্রথমে
স্থানান্তর, তারপর ঘোরাই।

প্রমাণপদ্ধতি আবার দেখি। সমতলের ঘূর্ণন-দল
সম্বন্ধে কিছু বীজগাণিতিক তথ্য দিয়ে শুরু করেছি।
এই দল দিয়ে $\onesphere$-কে তুল্যতা
শ্রেণিতে ভাগ করেছি। তারপর প্রত্যেক শ্রেণি
থেকে সদস্য নিতে চয়ন ব্যবহার করে
``কূটাভাস'' পেয়েছি।

বানাখ--তার্স্কিতেও ঠিক একই পদ্ধতি লাগে।
মূল পার্থক্য, তার জন্য প্রয়োজনীয়
বীজগাণিতিক তথ্যগুলি ভিতালির কূটাভাসের
চেয়ে অনেক জটিল। কিন্তু ওই তথ্যগুলির সঙ্গে
চয়নের সম্পর্ক নেই। দ্রুত সেগুলি সংক্ষেপ করি।

বানাখ--তার্স্কি প্রমাণে প্রথমে
\olref{disjointgroup}-এর অনুরূপ একটি
ফল প্রতিষ্ঠা করি: কমপক্ষে দুই উৎপাদকযুক্ত
% BN-SRC-473: not every free group is paradoxically decomposable;
% the nonabelian rank-at-least-two case is intended.
অ-আবেলীয় \emph{মুক্ত দল} চার ভাগে
ভাগ করা যায়; স্বজ্ঞায় ভাগগুলি ``সরিয়ে''
পুরো দলের দুটি প্রতিলিপি উদ্ধার করা যায়।\footnote{
\emph{চারটি} ভাগে কাজ হওয়ার ফলটি
\cite{Robinson1947}-এর। নতুন প্রমাণের জন্য
দেখো \citet[Theorem 5.2]{Wagon2016}।
দলের খণ্ডগুলি ``সরানো'' বলার রীতিতে আমরা
\citet[p.~3]{Weston2003}-কে অনুসরণ করছি।}
তারপর দেখাই, $\Real^3$-এর মূলবিন্দুকে
কেন্দ্র করে দুটি নির্দিষ্ট ঘূর্ণন দিয়ে
ঘূর্ণনের একটি মুক্ত দল $F$ উৎপন্ন হয়।\footnote{
দেখো \citet[Theorem 2.1]{Wagon2016}।}
(এখনও চয়ন লাগেনি।) গোলকের পৃষ্ঠের দুটি
বিন্দু তখনই ``একই ধরনের'', যখন $F$-এর
কোনো ঘূর্ণনে একটিকে অন্যটিতে নেওয়া যায়।
এরপর \emph{চয়ন ব্যবহার করে} প্রত্যেক
``একই ধরনের'' শ্রেণি থেকে ঠিক একটি
বিন্দু নিই। $F$-এর চার ভাগ গোলকের
পৃষ্ঠে \olref{pseudobanachtarski}-এর
মতো প্রয়োগ করে পৃষ্ঠটিকে চার ভাগে
ভাগ করি, যেগুলি ``সরিয়ে'' পৃষ্ঠের
দুটি প্রতিলিপি পাই। এভাবেই
\citep{Hausdorff1914} প্রতিষ্ঠিত হয়:
% BN-SRC-830 TARGET-CORRECTION-BEGIN
\par\noindent\textbf{পাঠ-সংশোধন।} বৃত্তে একটি অ-অভেদ ঘূর্ণন কোনও বিন্দু স্থির রাখে না। কিন্তু
ত্রিমাত্রিক গোলকের পৃষ্ঠে প্রত্যেক অ-অভেদ ঘূর্ণন তার অক্ষের
দুটি বিন্দু স্থির রাখে। তাই বিন্দুর একই image থেকে দুই
ঘূর্ণন অভিন্ন—আগের ধাপটি গোটা পৃষ্ঠে চলে না। ধরি $D$ হলো
$F$-এর কোনও অ-অভেদ ঘূর্ণনে স্থির থাকা পৃষ্ঠ-বিন্দুর সেট।
$F$ তালিকায়নযোগ্য এবং প্রত্যেক ঘূর্ণনের দুই স্থিরবিন্দু,
তাই $D$ তালিকায়নযোগ্য; conjugation-এ এটি $F$-অপরিবর্তিত।
পৃষ্ঠ থেকে $D$ বাদ দিলে দলটির ক্রিয়া মুক্ত, তাই সেখানেই
প্রতিনিধি নিয়ে দলের বিভাজন স্থানান্তর করা যায়।
ব্যতিক্রমগুলি ফেরানোর ধাপও দরকার। $D$-র কোনও বিন্দু দিয়ে
যায় না এমন একটি অক্ষ নিই। এই অক্ষে এমন ঘূর্ণন $\rho$
আছে যাতে $D,\rho D,\rho^2D,\ldots$ পরস্পর বিযুক্ত:
কোনও $d,e\in D$ এবং ধনাত্মক পূর্ণসংখ্যা $k$-র জন্য
$\rho^k d=e$ হওয়ার নিষিদ্ধ কোণ সসীমসংখ্যক, সব জোড়া ও
$k$ মিলে তালিকায়নযোগ্য কোণ বাদ যায়। বাকি কোণ থেকে একটি
নিই। ধরি $H=\bigcup_{n<\omega}\rho^nD$। $H$-তে
$\rho$ এবং বাকি পৃষ্ঠে অভেদ প্রয়োগে পুরো পৃষ্ঠ থেকে
$D$-বাদ দেওয়া পৃষ্ঠে দুটি খণ্ডের সমরূপ চিত্রণ হয়।
এই বিভাজন, মুক্ত-ক্রিয়ার বিভাজন এবং তার বিপরীত সংযোজন করে
খণ্ডগুলি প্রয়োজনে আরও ভাগ করলে পুরো পৃষ্ঠের সসীম বিভাজন
পাওয়া যায়। এই সংক্ষিপ্ত প্রমাণে গোটা পৃষ্ঠে ঠিক চার খণ্ড
দাবি করা হয়নি; খণ্ডসংখ্যার তীক্ষ্ণ ফলের অতিরিক্ত যুক্তি লাগে।
% BN-SRC-830 TARGET-CORRECTION-END

\begin{thm}[হাউসডর্ফের কূটাভাস ($\ZFC$-এ)]
যেকোনো গোলকের পৃষ্ঠকে সসীমসংখ্যক খণ্ডে
ভাগ করে, সেগুলি ঘুরিয়ে ও স্থানান্তর করে
আবার সাজালে সেই পৃষ্ঠের দুটি পৃথক
প্রতিলিপি গঠন করা যায়।
\end{thm}

পূর্ণ বানাখ--তার্স্কি উপপাদ্য পেতে
আরও কয়েকটি বীজগাণিতিক কৌশল লাগে:
সেখানে শুধু গোলকের পৃষ্ঠ নয়, ভেতরটিও
আছে। তবে এগুলি বীজগণিতের পিঠার
উপর সাজসজ্জার মতো। তাই ওয়েস্টন লেখেন:
\begin{quote}
  [\ldots] মুক্ত দল নিয়ে ফলটিই
  বানাখ--তার্স্কি কূটাভাসের প্রমাণের
  \emph{মূল ধাপ}। এই দিক থেকে কূটাভাসটি
  $\Real^3$ নিয়ে যতটা, তার চেয়ে বেশি
  [$\Real^3$-এর স্থানান্তর ও ঘূর্ণনের]
  দলের জটিলতা নিয়ে। \cite[p.~16]{Weston2003}
\end{quote}
অর্থাৎ \emph{সসীম} বিভাজন (বানাখ--তার্স্কি)
নাকি \emph{তালিকায়নযোগ্য অসীম} বিভাজন
(ভিতালি)—কোনটি সম্ভব, তা দ্বিমাত্রা বা
ত্রিমাত্রায় কাজের দলতাত্ত্বিক তথ্যের উপর নির্ভর করে।

স্বীকার করি, শেষ পর্যবেক্ষণটি
\olref[banach]{sec}-এর শেষের রসিকতা
একটু নষ্ট করে। ``Banach-Tarski''
দ্বিমাত্রিক বলে একে
``Banach-Tarski Banach-Tarski'' করতে
\emph{তালিকায়নযোগ্য অসীমসংখ্যক} খণ্ডে
ভাগ করতে হবে। রসিকতাটি ঠিক করতে
ত্রিমাত্রায় লিখতে হবে। পাঠকের অনুশীলন রইল।

শেষ মন্তব্য। \olref[banach]{sec}-এ
বলেছিলাম, গোলকের ``খণ্ডগুলি''
\emph{পরিমাপযোগ্য} হতে পারে না;
তারা ছবি আঁকা যায় না এমন ``অসীমভাবে
ছড়ানো'' সমাবেশ। \olref{pseudobanachtarski}
পেতে চয়নের যে ব্যবহার, সেখানেও তাই।
কথাটি ব্যাখ্যা করা উচিত।

আবার কিছু পটভূমি সংক্ষেপে বলি (এটি
\emph{শুধুই} রূপরেখা; \emph{পরিমাপ}
নিয়ে পাঠ্যবই দেখতে পারো)। $X$-এর
কোনো সেটের পরিমাপ দিতে হলে $X$-এর
কোনো ``$\sigma$-বীজগণিত''-এর প্রত্যেক
$E$-র জন্য $\mu(E)\in\Real$ মান দিতে
হয়। বিশদ এখানে জরুরি নয়; শুধু $\mu$-কে
তালিকায়নযোগ্য যোগধর্ম মানতে হবে:
বিযুক্ত সেটগুলির তালিকায়নযোগ্য সংযোগের
পরিমাপ তাদের পৃথক পরিমাপের যোগফল। অর্থাৎ
$X_n$-গুলি বিযুক্ত হলে
$\mu(\bigcup_{n<\omega}X_n)=
\sum_{n<\omega}\mu(X_n)$। কোনো সেট
``অপরিমাপযোগ্য'' মানে, উপযুক্তভাবে তার
পরিমাপ নির্ধারণ করা যায় না।
% BN-SRC-831 TARGET-CORRECTION-BEGIN
\par\noindent\textbf{পাঠ-সংশোধন।} অপরিমাপযোগ্যতা এখানে নির্দিষ্ট পরিমাপ ও তার সংজ্ঞাক্ষেত্রের
সাপেক্ষ; কোনও পরিমাপই দেওয়া যায় না, এমন দাবি নয়। যেমন
সব উপসেটে শূন্য পরিমাপ দেওয়া যায়, কিন্তু তা নিচের
$\mu(\onesphere)=1$ শর্ত মেটায় না। পরিমাপের মান সাধারণত
অঋণাত্মক প্রসারিত বাস্তব সংখ্যা; নিচের স্বাভাবিকীকৃত বৃত্তে
প্রত্যেক পরিমাপযোগ্য উপসেটের মান $0$ থেকে $1$-এর মধ্যে।
নিচের অনুসিদ্ধান্তে সংজ্ঞাক্ষেত্রটি ঘূর্ণন-অপরিবর্তিত
$\sigma$-বীজগণিত, আর ঘূর্ণনে পরিমাপ অক্ষত—এই শর্তগুলি
দরকার। একটি $\rho C$ পরিমাপযোগ্য হলে সব ঘূর্ণিত প্রতিলিপিও
তাই, এবং তাদের মান একই অঋণাত্মক $r$। বিযুক্ত গণনাযোগ্য
সংযোগের মান তখন $r=0$-তে শূন্য এবং $r>0$-তে অসীম,
যা বৃত্তের একক পরিমাপের সঙ্গে বিরোধ।
% BN-SRC-831 TARGET-CORRECTION-END
আগের
$\rotationsgroup$ নিয়ে এখন:

\begin{cor}[ভিতালি]
$\mu$ এমন পরিমাপ হোক যে
$\mu(\onesphere)=1$ এবং $X$ ও $Y$
সর্বসম হলে $\mu(X)=\mu(Y)$।
তাহলে প্রত্যেক $\rho\in\rotationsgroup$-র
জন্য $\funimage{\rho}{C}$ অপরিমাপযোগ্য।
\end{cor}

\begin{proof}
বিরোধের জন্য অন্যথা ধরি। অর্থাৎ কোনো
$\sigma\in\rotationsgroup$ ও
$r\in\Real$-র জন্য
$\mu(\funimage{\sigma}{C})=r$ ধরি।
% BN-SRC-475: rotations are indexed by rotationsgroup, not by C.
প্রত্যেক $\rho\in\rotationsgroup$-র জন্য
$\funimage{\rho}{C}$ ও
$\funimage{\sigma}{C}$ সর্বসম; তাই
প্রত্যেক $\rho\in\rotationsgroup$-র
জন্য $\mu(\funimage{\rho}{C})=r$।
\olref{vitalicover} ও
\olref{vitalinooverlap} অনুযায়ী
$\onesphere=\bigcup_{\rho\in\rotationsgroup}
\funimage{\rho}{C}$ জোড়ায় জোড়ায়
বিযুক্ত সেটের তালিকায়নযোগ্য সংযোগ।
তাই তালিকায়নযোগ্য যোগধর্মে
$\mu(\onesphere)=1$ হলো প্রত্যেক
$\funimage{\rho}{C}$-র পরিমাপের যোগফল:
\[
	1 = \mu(\onesphere) = \sum_{\rho \in \rotationsgroup}\mu(\funimage{\rho}{C}) = \sum_{\rho \in \rotationsgroup}r
\]
কিন্তু $r=0$ হলে
$\sum_{\rho\in\rotationsgroup}r=0$,
আর $r>0$ হলে
$\sum_{\rho\in\rotationsgroup}r=\infty$।
পরিমাপ ঋণাত্মক নয়; তাই উভয় ক্ষেত্রেই
১-এর সঙ্গে বিরোধ।
\end{proof}

\end{document}
